%% P27 --- Wnioskowanie statystyczne dla inżyniera
%%
%% Bliźniak: programs/en/P27-inference.tex. Ramka w ramkę, makro w makro,
%% liczba w liczbę; tools/parity.py porównuje oba pliki. Każda liczba
%% pochodzi z code/p27_inference.py.
%%
%% CO WYDALI SĄSIEDZI --- przeczytane, zanim powstała pierwsza linijka.
%% Pełna lista jest w nagłówku skryptu. W skrócie: DWA Z SIEDMIU PUNKTÓW
%% BRIEFU SĄ JUŻ NA STRONIE:
%%   P25 sekcja 4 ma błąd standardowy proporcji, regułę czterech razy tyle
%%       pozycji ORAZ RACHUNEK MOCY -- 12293 pozycje na model, żeby zobaczyć
%%       jeden punkt. Ten program z tego korzysta i niczego nie powtarza. To,
%%       czego P25 nie mógł zrobić, to ten sam rachunek dla modeli, które
%%       widziały TE SAME prompty; jego własna sekcja 1 daje wyraz z
%%       kowariancją. Sekcja 3 to liczba z P25 zrobiona porządnie i związana
%%       z nią bramką.
%%   P25 sekcja 1 NAZYWA też bootstrap po wierszach jednego użytkownika jako
%%       tryb awarii, zanim książka powie, czym bootstrap jest. Sekcja 2 to
%%       spłaca.
%%   P26 sekcja 1 ma obciążenie, wariancję i błąd średniokwadratowy i celowo
%%       zatrzymuje się przed pytaniem, czy różnica jest prawdziwa.
%%   P23 odsyła "co p-wartość mówi, a czego nie" TUTAJ z nazwy i ma odwrócenie
%%       warunkowania, którym błąd p-wartości jest. Sekcja 4 to twierdzenie z
%%       P23 przeczytane na innej parze zdarzeń.
%%   P12 ma symbol Newtona. Test dokładny z sekcji 3 to math.comb na liczbach
%%       całkowitych i nic więcej, dzięki czemu p-wartość pojawia się jako
%%       liczba rzutów monetą, zanim pojawi się jako słowo.
%%   F10 ma licznik i mianownik. Sekcja 1 to ta lekcja zastosowana do miejsca
%%       dziesiętnego w wyniku benchmarku.
%%
%% CELOWO NIEWYDANE TUTAJ:
%%   P28 ma rozkład a posteriori, sprzężoność, przedział wiarygodności i
%%       "prawdopodobieństwo, że B jest lepszy". Sekcja 4 mówi wprost, że
%%       p-wartość tym nie jest, i wskazuje P28; nie wolno jej tu liczyć.
%%   P34 ma projektowanie ewaluacji od początku do końca.

\program[Wnioskowanie statystyczne]{Wnioskowanie statystyczne dla inżyniera}
\label{prog:P27}
\index{wnioskowanie!statystyczne}
\index{ewaluacja!czy różnica jest prawdziwa}

\begin{outcomes}
\outcome{1--6}{powiedzieć, ile warta jest podawana różnica w benchmarku
  \dash{} w pozycjach, a nie w punktach}
\outcome{7--14}{zbudować przedział ufności dla metryki, która nie ma wzoru, i
  nazwać jednostkę, którą trzeba losować}
\outcome{15--22}{zaplanować i przeprowadzić porównanie dwóch modeli
  ewaluowanych na tych samych promptach}
\outcome{23--30}{powiedzieć, czym p-wartość jest, a czym nie jest}
\outcome{31--36}{powiedzieć, co czterdzieści porównań w jednym rankingu robi z
  tym wszystkim}
\outcome{37--39}{zastosować całość do liczby, którą ktoś opublikował}
\end{outcomes}

\begin{quiz}
\item Model ma według raportu $\val{p27.quoted.hi}$ procent na
  $\val{p27.items}$ pozycjach ewaluacyjnych. Co jest nie tak z tym zdaniem,
  jeszcze przed jakąkolwiek statystyką? \teachesat{1--2}
  \answerto{$\val{p27.quoted.hi}$ procent z $\val{p27.items}$ to
  $\val{p27.k.hi}$ pozycji, a modelowi nie da się zaliczyć
  $\val{p27.k.hi}$ pozycji. Na $\val{p27.items}$ pozycjach osiągalne wyniki
  dzieli $\val{p27.grid.pts}$ punktu, więc miejsce dziesiętne jest gęstsze
  niż siatka.}
\item Dwa modele ewaluowano na tych samych $\val{p27.items}$ promptach i jeden
  ma $\val{p27.gap.pts}$ punktu więcej. Z ilu pozycji składa się ta różnica?
  \teachesat{3--4}
  \answerto{$\val{p27.gap.items}$. Cała opublikowana różnica to jedna pozycja
  ewaluacyjna, która poszła w drugą stronę.}
\item Masz jeden zbiór ewaluacyjny i metrykę bez wzoru na błąd standardowy.
  Wskaż sposób, żeby mimo to dostać przedział. \teachesat{7--8}
  \answerto{Bootstrap: losuj pozycje ze zwracaniem, licz metrykę na każdym
  takim losowaniu i weź środek otrzymanych odpowiedzi.}
\item Dwa modele ewaluowano na tych samych promptach, a nie na osobnych
  próbach. Czy to czyni porównanie trudniejszym, czy łatwiejszym?
  \teachesat{16--17}
  \answerto{Łatwiejszym, i zwykle znacznie. Oba wyniki są skorelowane, a
  wariancja różnicy odejmuje podwojoną kowariancję, więc potrzebna liczba
  pozycji spada o czynnik $1-\rho$.}
\item Porównanie zwraca $p = \num{0.03}$. Jakie jest prawdopodobieństwo, że
  oba modele są równie dobre? \teachesat{24--25}
  \answerto{P-wartość tego nie mówi. To prawdopodobieństwo wyniku równie
  skrajnego \emph{przy założeniu}, że są równie dobre, a czytanie tego wstecz
  to odwrócenie, przed którym ostrzega Program~\ref{prog:P23}.}
\end{quiz}

% ============================================================
\section{Ile wart jest jeden punkt}
\label{sec:P27-one-point}
\index{ewaluacja!ziarnistość wyniku}
\index{pomiar!czytanie miejsca dziesiętnego wobec mianownika}

\begin{fr}
Oto twierdzenie z gatunku, który czytelnik tej książki spotyka co tydzień, i
to na nie cały program odpowiada.

\begin{center}
\emph{Model B dostał $\val{p27.quoted.hi}$, a model A
$\val{p27.quoted.lo}$ na $\val{p27.items}$ pozycjach ewaluacyjnych.}
\end{center}

Rozstrzygnięcie, czy to prawdziwa różnica, wymaga narzędzi. Ale jedna rzecz
nie wymaga żadnej statystyki i warto ją znaleźć jako pierwszą.

Wynik to pewna liczba pozycji zaliczonych z $\val{p27.items}$. Czy model może
dostać $\val{p27.quoted.hi}$ procent na $\val{p27.items}$ pozycjach?
\dotline
\nextframe
\end{fr}

\begin{fr}
\ans{Nie}

$\val{p27.quoted.hi}$ procent z $\val{p27.items}$ to $\val{p27.k.hi}$
pozycji, a nikt nigdy nie zaliczył $\val{p27.k.hi}$ pozycji. Osiągalne wyniki
to $\val{p27.near.lo}$ i $\val{p27.near.hi}$, i nie ma między nimi niczego,
ponieważ \textbf{wynik $k$ zaliczonych z $n$ leży na siatce o oczku
$100/n$ punktu} \dash{} tutaj $\val{p27.grid.pts}$ punktu.
Druga liczba nie przechodzi tego samego testu.

Zdanie nie jest więc wewnętrznie spójne, a przydatna reakcja pochodzi z
Programu~\ref{prog:F10}: \textbf{zapytaj, co jest mianownikiem.} Albo
ewaluacja nie miała $\val{p27.items}$ pozycji, albo wynik nie jest ich
proporcją \dash{} na przykład jest średnią po podbenchmarkach różnej
wielkości, a wtedy $\val{p27.items}$ też nie jest liczbą, z której należy
liczyć przedział.

Żadne z tych odczytań nie jest niezwykłe i żadne nie stoi w rankingu.
Załóżmy, że benchmark chce podawać dokładność z jednym miejscem po przecinku.
Ilu pozycji to wymaga, zanim to miejsce w ogóle zaokrągli się poprawnie?
\dotline
\nextframe
\end{fr}

\begin{fr}
\ans{Ponad tysiąca: $\val{p27.decimal.floor}$}

Siatka musi być gęstsza niż dziesiąta punktu, więc $1/n$ musi być
poniżej $\num{0.001}$.

To stwierdzenie o arytmetyce, a za nim stoi drugie zastrzeżenie, dużo
większe. Na $\val{p27.decimal.floor}$ pozycjach \emph{siatka} jest już dość
gęsta, żeby wydrukować miejsce dziesiętne, a przedział z
Programu~\ref{prog:P25} na tylu pozycjach to wciąż jakieś trzy punkty w każdą
stronę. Pierwsze zastrzeżenie mówi, że cyfry nie da się policzyć; drugie, że
nic by nie znaczyła, gdyby się dało.

Przyjmijmy mimo to twierdzenie na wiarę i zadajmy pytanie, które rozstrzyga,
czy cokolwiek z tego ma znaczenie. Różnica to $\val{p27.gap.pts}$ punktu. Ile
to w pozycjach?
\dotline
\nextframe
\end{fr}

\begin{fr}
\ans{$\val{p27.gap.items}$}

\textbf{Jedna pozycja.} Na $\val{p27.items}$ pozycjach różnica
$\val{p27.gap.pts}$ punktu to jedno pytanie, na które B odpowiedział, a A nie,
i to jest cała zaobserwowana różnica między dwoma modelami.

Warto przy tym usiąść, bo opublikowane zdanie tego nie mówi. Procent ukrywa
własny mianownik, a dwie liczby oddalone o kilka dziesiątych czytają się jak
zmierzone rozróżnienie, a nie jak jedno pytanie o stolicę czegoś.

Program~\ref{prog:P25} wycenił szum wokół takiej liczby i ten program tego nie
wyprowadza od nowa: przy $\val{p27.quoted.hi}$ procent na $\val{p27.items}$
pozycjach przedział to około $\pm\val{p27.hw.200}$ punktu. Zestaw to z
różnicą $\val{p27.gap.pts}$ punktu.
\dotline
\nextframe
\end{fr}

\begin{fr}
\ans{Przedział jest ponad $\val{p27.hw.ratio.floor}$ razy większy}

Co nie jest sytuacją graniczną.

\mermaidfig{p27-one-item}{Opublikowana różnica w jednostce, którą ewaluacja
  naprawdę ma}{różnica w pozycjach}

To wszystko, co ustala sekcja 1, i nie potrzeba do tego żadnego wnioskowania
\dash{} dwa dzielenia i pierwiastek z Programu~\ref{prog:P25}. Reszta programu
dotyczy przypadków, w których odpowiedź nie jest tak oczywista, oraz jednego
odruchu, który zaraz pogorszy następne porównanie czytelnika, zamiast je
poprawić.
\end{fr}

\begin{fr}
\begin{trapbox}
\textbf{\enquote{Przedziały się nakładają, więc różnica nie jest prawdziwa.}}

To najczęstszy wniosek wyciągany z pary słupków błędu i nie jest to żaden
test. Dwa przedziały mogą nakładać się znacznie, podczas gdy różnica między
wielkościami spokojnie odróżnia się od zera, i mogą się nie nakładać, gdy tak
nie jest.

Powód należy do sekcji 3 i warto go podać wcześnie: przedział na modelu A i
przedział na modelu B dotyczą każdy \emph{jednego modelu}, a pytanie dotyczy
\emph{różnicy}, która jest trzecią wielkością z własnym rozrzutem.
Program~\ref{prog:P25} pokazał już, że różnica dwóch niezależnych ocen niesie
sumę ich wariancji. Gdy oba mierzono na tych samych promptach, niesie coś
jeszcze mniejszego, a porównanie dwóch osobnych przedziałów nie widzi ani
jednego, ani drugiego.

Nawyk wart nabycia: \textbf{buduj przedział na wielkości, o którą naprawdę
pytasz}, czyli tutaj na różnicy, a nie na dwóch wynikach.
\end{trapbox}

Wiesz już, ile warta jest różnica na tym zbiorze ewaluacyjnym. Następne
pytanie brzmi, co robić, gdy metryka nie jest dokładnością i nie ma
pierwiastka, po który można sięgnąć.
\end{fr}

% ============================================================
\section{Przedział bez wzoru}
\label{sec:P27-bootstrap}
\index{bootstrap}
\index{ewaluacja!przedział ufności}

\begin{fr}
Przedział z Programu~\ref{prog:P25} działa, bo dokładność jest średnią $n$
niezależnych trafień lub pudeł, a wariancja jednego z nich to $p(1-p)$.
Prawie nic innego, co raportujesz, nie jest tak wygodne. Mediana opóźnienia,
odsetek powodzeń przy $k$ próbach, średnia makro po klasach, odsetek wygranych
oceniany przez inny model: żadne z nich nie ma błędu standardowego w
podręczniku leżącym na twoim biurku.

Załóżmy, że pieniądze nie grają roli i możesz przeprowadzić całą ewaluację od
nowa, tysiąc razy. Co byś zrobił z tysiącem odpowiedzi?
\dotline
\nextframe
\end{fr}

\begin{fr}
\begin{ansblock}
Zobaczył, jak bardzo się różniły \dash{} wziął środkowe $95$ procent z nich i
to podał.
\end{ansblock}

To definicja tego, czym przedział próbuje być, i nie wymaga wzoru dla żadnej
metryki. Przeszkodą jest to, że nie da się tego zrobić: masz jeden zbiór
ewaluacyjny i jedną odpowiedź.

Ruch bootstrapu to jedno zdanie. \textbf{Próbka, którą masz, jest najlepszym
dostępnym opisem populacji, z której pochodzi, więc losuj nowe próbki z
niej.} Wylosuj $n$ pozycji ze swoich $n$ pozycji \emph{ze zwracaniem}, policz
metrykę i powtórz.

Weź przypadek, w którym wzór jednak masz, żeby dało się porównać. Twoja
ewaluacja to $\val{p27.boot.n}$ pozycji, a model zaliczył $\val{p27.boot.k}$ z
nich. Losujesz $\val{p27.boot.n}$ pozycji ze zwracaniem i liczysz poprawne.
Jaki to rozkład?
\dotline
\nextframe
\end{fr}

\begin{fr}
\ans{Dwumianowy}

Dokładnie $\mathrm{Binomial}(\val{p27.boot.n}, \val{p27.boot.k}/\val{p27.boot.n})$,
bo każde losowanie trafia na poprawną pozycję z prawdopodobieństwem
$\val{p27.boot.k}/\val{p27.boot.n}$, a losowania są niezależne.

Na proporcji \textbf{rozkład bootstrapowy nie jest więc przybliżeniem
czegokolwiek} \dash{} jest rozkładem dwumianowym i nie trzeba żadnej symulacji,
żeby go znaleźć. Przedział to kwantyle tego rozkładu, otrzymane przez
zsumowanie
symbolu Newtona z Programu~\ref{prog:P12}:

\begin{center}
\begin{tabular}{lrr}
\toprule
& dolny & górny \\
\midrule
bootstrap & $\val{p27.boot.lo}$ & $\val{p27.boot.hi}$ \\
wzór z Programu~\ref{prog:P25} & $\val{p27.closed.lo}$ &
  $\val{p27.closed.hi}$ \\
\bottomrule
\end{tabular}
\end{center}

Oba zgadzają się co do $\val{p27.boot.agree}$ punktu, czyli jednej pozycji
\dash{} najdrobniejszej rzeczy, jaką którykolwiek z nich rozróżnia, według
sekcji 1. Skrypt sprawdza właśnie to, a nie miejsce dziesiętne, bo zgodność
dokładniejsza niż siatka byłaby twierdzeniem, którego żadna z metod nie
uniesie.
\end{fr}

\begin{fr}
Zgodność uspokaja i nie o nią chodzi. Gdyby bootstrap wyłącznie odtwarzał
wzór, który już masz, byłby ciekawostką.

Oto przypadek, dla którego istnieje. Dwadzieścia zmierzonych opóźnień w
milisekundach:

\begin{center}
\code{112 118 121 125 129 131 134 138 141 147}\\
\code{152 158 163 171 180 195 214 248 310 502}
\end{center}

Mediana to $\val{p27.med.obs}$ ms. Zanim pójdziesz dalej, wypisz wzór na błąd
standardowy mediany.
\dotline
\nextframe
\end{fr}

\begin{fr}
\begin{ansblock}
Żadnego znanego nie ma, a ten, który istnieje, potrzebuje gęstości rozkładu w
medianie \dash{} czyli wielkości, którą trzeba by oszacować z tej samej próbki,
przyjmując założenie niesprawdzalne.
\end{ansblock}

Bootstrap nie potrzebuje niczego z tego. Wylosuj dwadzieścia opóźnień z tych
dwudziestu ze zwracaniem, weź medianę i zrób tak $\val{p27.med.trials}$ razy:

\[ \val{p27.med.obs}\ \text{ms}, \quad
   \text{przedział } \intcc{\val{p27.med.lo}}{\val{p27.med.hi}} \]

Zwróć uwagę, czym są końce. Mediana wylosowanej próbki jest zawsze wartością z
pierwotnej listy albo środkiem dwóch takich wartości, więc \textbf{końce
przedziału to pomiary, które ktoś wykonał}, a nie wyjścia dopasowanego modelu.
To też powód, dla którego odpowiedź jest stabilna: inny generator liczb
losowych wybiera inne próbki i trafia na te same wartości z danych.

Procedura nie zmieniła się ani trochę między tymi dwoma przypadkami.
Podmieniono metrykę i nic poza tym, i to jest własność warta posiadania.
\end{fr}

\begin{fr}
\begin{rigourbox}
\textbf{Nie dowodzimy tutaj: dlaczego bootstrap działa.}

Uzasadnienie brzmi tak, że empiryczny rozkład próbki zbiega do rozkładu, z
którego ją wylosowano, więc zachowanie statystyki przy losowaniu z próbki
zbiega do jej zachowania przy losowaniu z populacji. Twierdzenie, jego
założenia i przypadki, w których zawodzi, to kurs magisterski, a oryginał to
Efron, 1979.

Przy tym biurku liczy się kształt awarii, a nie dowód. \textbf{Bootstrap nie
wymyśli pozycji, których nie ewaluowałeś.} Jeśli twoje dwieście promptów nie
zawiera zadań arytmetycznych, to żadne wylosowanie z nich ich nie zawiera, a
przedział jest uczciwy wobec błędu losowania i milczy o tym, co naprawdę cię
ugryzie. Najgorzej radzi sobie też ze skrajnościami \dash{} z maksimum, z
$99$-tym centylem \dash{} bo próbka może sięgnąć tylko po wartości już obecne.
\end{rigourbox}

Jeszcze jedna decyzja, zanim technika nadaje się do użytku, i
Program~\ref{prog:P25} nazwał ją jako tryb awarii, zanim ta książka
powiedziała, czym bootstrap jest. Twoja ewaluacja ma $\val{p27.clust.n}$
wierszy, ale pochodzą one od $\val{p27.clust.users}$ użytkowników po
$\val{p27.clust.rows}$ promptów każdy. Co losujesz?
\dotline
\nextframe
\end{fr}

\begin{fr}
\ans{Użytkownika, a nie wiersz}

A pomyłka tutaj nie jest błędem zaokrąglenia. Powodem jest własna tożsamość z
Programu~\ref{prog:P25}: wariancja sumy rzeczy \emph{niezależnych} rośnie jak
ich liczba, a wariancja sumy rzeczy \emph{identycznych} jak kwadrat tej
liczby. Wiersze jednego użytkownika leżą gdzieś pomiędzy, a w skrajności są
jedną obserwacją noszącą $\val{p27.clust.rows}$ etykiet.

Bootstrap po wierszach uważa, że ma $\val{p27.clust.n}$ niezależnych
obserwacji. Jeśli wiersze wewnątrz użytkownika zgadzają się całkowicie, ma
$\val{p27.clust.users}$, a przedział, który podaje, jest za wąski o czynnik
$\val{p27.clust.factor}$ \dash{} pierwiastek z $\val{p27.clust.rows}$.

\mermaidfig{p27-what-to-resample}{Jednostka losowania to decyzja, której nic
  nie sprawdza}{jednostka}

Żadna biblioteka ci tego nie powie. Funkcja przyjmuje tablicę, a tablica
$\val{p27.clust.n}$ wierszy wygląda dokładnie jak $\val{p27.clust.n}$
obserwacji.
\end{fr}

\begin{fr}
\begin{trapbox}
\textbf{\enquote{Zrób więcej losowań, a przedział się zawęzi.}}

Nie zawęzi się, a przekonanie warto zabić, bo losowania są tanie i robienie
ich więcej sprawia wrażenie działania.

W grze są dwa różne błędy. \emph{Szerokość} przedziału wynika z tego, ile
informacji zawiera zbiór ewaluacyjny, czyli z $n$ pozycji i z niczego innego.
Liczba losowań rządzi wyłącznie tym, jak dokładnie namierzyłeś końce tego
przedziału \dash{} błędem Monte Carlo, czyli tym, co Program~\ref{prog:P25} zapisuje jako
$1/\sqrt{B}$ przy $B$ losowaniach i co możesz zbić dowolnie nisko za cenę
arytmetyki.

Więcej losowań kupuje więc dokładniej podany przedział \textbf{tej samej
szerokości}. Węższy kosztuje pozycje, po czterokroć za połowę z
Programu~\ref{prog:P25}.
\end{trapbox}

Umiesz już nałożyć przedział na cokolwiek, co potrafisz policzyć. Zostaje
pytanie, od którego program się zaczął, i okazuje się, że zadanie go wprost
jest znacznie tańsze niż nakładanie przedziału osobno na każdy model.
\end{fr}

% ============================================================
\section{Widziały te same prompty}
\label{sec:P27-paired}
\index{ewaluacja!porównanie sparowane}
\index{pomiar!planowanie porównania}

\begin{fr}
Program~\ref{prog:P25} odpowiedział na pytanie o rozmiar i odpowiedź była
zniechęcająca: rozstrzygnięcie różnicy jednego punktu między dwoma modelami
kosztuje $\val{p27.n.rho00}$ pozycji ewaluacyjnych \emph{na model}, czyli
czyjś tygodniowy budżet i więcej promptów, niż zawiera większość benchmarków.

Ta liczba jest poprawna i zawiera założenie, które prawie nigdy nie jest
prawdziwe o tym, co naprawdę robisz. Przeczytaj zdanie jeszcze raz i nazwij
je.
\dotline
\nextframe
\end{fr}

\begin{fr}
\begin{ansblock}
Że obie ewaluacje są od siebie niezależne.
\end{ansblock}

Nie są. Oba modele puściłeś na \textbf{tych samych promptach}, bo to jedyny
rozsądny sposób ich porównania i bo zebranie drugiego zbioru ewaluacyjnego
kosztowałoby dwa razy tyle.

Sekcja 1 Programu~\ref{prog:P25} ma tożsamość, która to naprawia, i niosła już
wyraz, którego tu potrzeba:

\[ \Var(X - Y) \;=\; \Var(X) + \Var(Y) - 2\Cov(X, Y) \]

P25 użył przypadku niezależnego, w którym kowariancja jest zerem, a wariancje
się dodają \dash{} stamtąd wziął się jego czynnik dwa. Tutaj kowariancja nie
jest zerem. Powiedz, co to robi z liczbą pozycji, i powiedz to jako wzór, a
nie jako kierunek.
\dotline
\nextframe
\end{fr}

\begin{fr}
\begin{ansblock}
Mnoży ją przez $1 - \rho$, gdzie $\rho$ jest korelacją wyników obu modeli na
poszczególnych pozycjach.
\end{ansblock}

Oba modele leżą blisko tej samej dokładności $p$, więc każdy wskaźnik na
pozycji ma wariancję $p(1-p)$, a ich kowariancja to $\rho\,p(1-p)$. Różnica
niesie wtedy $2p(1-p)(1-\rho)$ tam, gdzie przypadek niezależny niósł
$2p(1-p)$, i cały rachunek z Programu~\ref{prog:P25} przechodzi bez zmian z
tym jednym czynnikiem doczepionym:

\begin{center}
\begin{tabular}{lrr}
\toprule
$\rho$ & pozycji na model & wobec $\rho = 0$ \\
\midrule
$0$ & $\val{p27.n.rho00}$ & $1\times$ \\
$\num{0.5}$ & $\val{p27.n.rho50}$ & $\tfrac{1}{2}\times$ \\
$\num{0.8}$ & $\val{p27.n.rho80}$ & $\tfrac{1}{5}\times$ \\
$\num{0.9}$ & $\val{p27.n.rho90}$ & $\tfrac{1}{10}\times$ \\
\bottomrule
\end{tabular}
\end{center}

Skrypt sprawdza czynnik $1-\rho$ przy czterech dokładnościach i trzech
wielkościach różnicy, a nie cztery liczby z tabeli, bo czynnik jest wynikiem,
a liczby są jedną jego komórką. A górny wiersz to zatwierdzona liczba samego
Programu~\ref{prog:P25}: wzór dla par napisano tak, żeby ją zawierał, i build
pada, gdyby kiedykolwiek przestały się zgadzać.
\end{fr}

\begin{fr}
Pytanie brzmi więc, jak duże jest $\rho$ w praktyce, a odpowiedź brzmi, że dla
dwóch modeli wartych porównania jest zwykle duże. Oba wytrenowano na
większości internetu, oba zaliczają łatwe pozycje i oba wykładają się na
pozycji, której pytanie jest niejednoznaczne. Niezgoda jest wyjątkiem, i
dokładnie to czyni porównanie tanim.

\mermaidfig{p27-same-prompts}{Skąd bierze się oszczędność w porównaniu
  sparowanym}{parowanie odejmuje}

To zmiana rzędu wielkości w koszcie eksperymentu, z jednej linijki algebry
samego Programu~\ref{prog:P25}. To także powód, dla którego rada
\enquote{ewaluuj oba modele na tych samych promptach} nie jest wyłącznie
schludnością.

Teraz zaostrzmy to. Z twoich $\val{p27.items}$ pozycji część oba modele
zaliczyły, część oba oblały, a na części dokładnie jeden zaliczył. Która z
tych trzech grup niesie informację o \emph{różnicy}?
\dotline
\nextframe
\end{fr}

\begin{fr}
\ans{Tylko trzecia}

Pozycja, na którą oba modele odpowiedziały poprawnie, wnosi tyle samo do obu
wyników i skraca się z różnicy dokładnie. Tak samo pozycja, którą oba oblały.
Tylko pozycje \textbf{niezgodne} \dash{} te, na których dokładnie jeden model
miał rację \dash{} mówią cokolwiek o tym, który model jest lepszy.

Załóżmy, że $\val{p27.disc}$ z $\val{p27.items}$ jest niezgodnych, i tak
właśnie wygląda wysokie $\rho$: oba modele zgodziły się na
$\val{p27.concord}$ pozycjach, więc porównanie opiera się na
$\val{p27.to.b}$ pozycjach dla B wobec $\val{p27.to.a}$ dla A. Sekcja 1
ustaliła, że B prowadzi o $\val{p27.net}$ pozycję, i zauważ, że nieparzysta
liczba pozycji niezgodnych jest tym, czego przewaga o jeden wymaga: liczba i
przewaga mają tę samą parzystość, bo przewaga jest różnicą dwóch liczb, które
sumują się do tej liczby.

Warto powiedzieć wprost, co znaczy tu \enquote{równie dobre}, bo na tym
opiera się dokładność tego, co dalej. To nie twierdzenie, że oba modele są tym
samym modelem. To słabsze twierdzenie, że na pozycji, na której się nie
zgadzają, każdy z nich mógł być tym, który ją zaliczył \dash{} więc każda
pozycja niezgodna jest rzutem monetą o to, który akurat nim był.

Teraz pytanie, i nie potrzeba do niego słownictwa statystycznego. Gdyby oba
modele były równie dobre w tym sensie, pozycje niezgodne rozkładałyby się na
jeden lub drugi jak uczciwe rzuty monetą. Rzucasz $\val{p27.disc}$ monetami.
Jak bardzo powinno cię zaskoczyć, że wypadły $\val{p27.net}$ od równego
podziału?
\dotline
\nextframe
\end{fr}

\begin{fr}
\ans{Ani trochę nie może cię zaskoczyć: prawdopodobieństwo wynosi
$\val{p27.p.exact}$}

\transcript{p27-odd-count}

\textbf{Dokładnie jeden}, i nie jest to zaokrąglenie. $\val{p27.disc}$ monet
nie może wypaść po równo, więc jedna strona zawsze prowadzi, a najmniejsza
dostępna przewaga to $\val{p27.net}$ \dash{} czyli ta, którą
zaobserwowano. Wynik co najmniej tak przechylony jest zatem pewny, a test nie
odróżnia tego wyniku od żadnego innego.

Powiedzieć da się to, który to wynik: przewaga dokładnie
$\val{p27.net}$, w jedną albo w drugą stronę, jest pojedynczym najbardziej
prawdopodobnym rezultatem i zdarza się w $\val{p27.p.one.pct}$ procentach
przypadków. \textbf{Opublikowana różnica między modelami jest najmniej
informatywnym wynikiem, na jaki arytmetyka pozwala.}

To symbol Newtona z Programu~\ref{prog:P12} na liczbach całkowitych i nic
więcej: żadnego przybliżenia normalnego, żadnej wariancji, żadnego
pierwiastka. Jest dokładny i identyczny na każdej maszynie.

Opublikowana różnica między modelami jest więc tym, czego spodziewałbyś się,
gdyby oba modele były takie same.

Co nasuwa pytanie w drugą stronę, a zadanie go nic nie kosztuje, bo to ta sama
suma policzona w górę zamiast raz. Zamiast sprawdzać, czy twój wynik
przekracza poprzeczkę, ustal, gdzie ta poprzeczka jest.

Jak duża musiałaby być różnica, żeby przestała być zwyczajna?
\dotline
\nextframe
\end{fr}

\begin{fr}
\ans{$\val{p27.net.needed}$ pozycji, czyli $\val{p27.net.needed.pts}$ punktu}

Znalezione przez szukanie w górę od $\val{p27.net}$, a nie zacytowane: na
$\val{p27.disc}$ pozycjach niezgodnych przewaga musi sięgnąć
$\val{p27.net.needed}$, zanim uczciwa moneta produkowałaby ją rzadziej niż raz
na dwadzieścia.

\textbf{Na $\val{p27.items}$ pozycjach, przy tych dwóch modelach, różnica
poniżej $\val{p27.net.needed.pts}$ punktu nie jest wynikiem.} To liczba, którą
czytelnik może zabrać na spotkanie, i jest trzynaście razy większa od tej,
którą opublikowano.

Dwa uczciwe zastrzeżenia i oba wzmacniają argument, zamiast go osłabiać. Próg
raz na dwadzieścia jest konwencją i następna sekcja mówi, co on znaczy, a
czego nie. Oraz $\val{p27.disc}$ pozycji niezgodnych zostało założone, a nie
zmierzone \dash{} to własność twoich dwóch modeli, nie drukuje jej żaden
harness ewaluacyjny i jest to najbardziej przydatna liczba, jaką mógłbyś do
niego dodać.
\end{fr}

\begin{fr}
\begin{aibox}
\textbf{Dwie linijki, które twój harness mógłby drukować i nie drukuje.}

Każdy raport ewaluacyjny w tej dziedzinie drukuje dwie dokładności. Prawie
żaden nie drukuje \textbf{liczby pozycji niezgodnych} i \textbf{przewagi}, a
te dwie liczby całkowite są całą treścią porównania \dash{} dokładności są ich
stratnym podsumowaniem. Mając je, powyższa arytmetyka to jedno wywołanie
\code{math.comb} i nie potrzeba żadnej biblioteki.

Koszt ich dodania to para liczników w pętli, która i tak porównuje oba modele
pozycja po pozycji. Koszt ich niedodania jest taki, że porównanie raportuje
się w postaci, z której nie da się go sprawdzić.
\end{aibox}
\end{fr}

% ============================================================
\section{Co mówi p-wartość}
\label{sec:P27-p-value}
\index{p-wartość}
\index{wnioskowanie!czego p-wartość nie mówi}

\begin{fr}
Liczba z poprzedniej sekcji ma nazwę, a ty policzyłeś ją już przed poznaniem
słowa, i to jest właściwa kolejność \dash{} w dodatku przy tej jednej
wartości, przy której najtrudniej pomylić się co do jej znaczenia.

To \textbf{p-wartość}, a oto definicja podana na tyle płasko, by dało się ją
zestawić z tym, co zrobiłeś:

\begin{center}
prawdopodobieństwo wyniku co najmniej tak skrajnego,\\
\emph{gdyby} oba modele były równie dobre
\end{center}

Każdy człon coś dźwiga. \enquote{Co najmniej tak skrajnego} to powód, dla
którego suma biegła po każdym wyniku równie przechylonym co ten, który
zobaczyłeś, a nie po nim samym. A \emph{gdyby} jest warunkiem, co czyni
całość prawdopodobieństwem warunkowym \dash{} tematem
Programu~\ref{prog:P23}.

Zapisz to jako takie, w notacji tamtego programu, a potem powiedz, co jest po
której stronie kreski.
\dotline
\nextframe
\end{fr}

\begin{fr}
\begin{ansblock}
$\Prob(\text{wynik tak skrajny} \mid \text{modele są równie dobre})$
\dash{} dane są zdarzeniem, a twierdzenie o modelach jest warunkiem.
\end{ansblock}

Czyli odwrotnie, niż to każdy czyta.

Załóżmy, że porównanie wraca z $p = \num{0.03}$. Czy jest trzyprocentowa
szansa, że oba modele są równie dobre?
\dotline
\nextframe
\end{fr}

\begin{fr}
\ans{Nie, a jedno z drugim nie wiąże się bez trzeciej liczby}

\begin{trapbox}
\textbf{\enquote{$p = \num{0.03}$, więc jest $3$-procentowa szansa, że to
przypadek.}}

To czyta $\Prob(A \mid B)$ jako $\Prob(B \mid A)$, czyli dokładnie
odwrócenie, któremu Program~\ref{prog:P23} poświęca sekcję, i jest to ten sam
błąd co przy częstości bazowej. P-wartość warunkuje na tym, że modele są
równe; zdanie warunkuje na danych. Własne twierdzenie
Programu~\ref{prog:P23} mówi, że jednego nie zamienisz w drugie bez rozkładu a
priori, a p-wartość nie przychodzi z żadnym.

Rozumowanie, które to produkuje, nie jest głupie, i dlatego warto je nazwać, a
nie skarcić. \enquote{To byłoby mało prawdopodobne, gdyby nic się nie działo,
więc coś się dzieje} jest poprawne na tyle, na ile sięga, i tak właśnie czyta
się test medyczny. Brakuje mu \dash{} i p-wartość tego nie dostarcza \dash{}
informacji, jak często coś się \emph{dzieje}, czyli częstości bazowej z
przykładu w Programie~\ref{prog:P23}.
\end{trapbox}
\end{fr}

\begin{fr}
Dajmy jej więc tę częstość i zobaczmy, czym staje się liczba. Maszyneria
Programu~\ref{prog:P23} stosuje się bez zmian: ustaw po jednej stronie
porównania, które ktoś przeprowadza w ciągu roku, i zapytaj, jaka część tych,
które wracają jako \enquote{istotne}, jest prawdziwa.

Zarzut wobec wstawiania tam liczby jest taki, że to zgadywanie, i warto na
niego odpowiedzieć, a nie go przyjąć. To zgadywanie, które i tak robisz w
domyśle za każdym razem, gdy czytasz wynik, i takie, które twój własny zespół
umie oszacować: ile ze zmian proponowanych w zeszłym roku faktycznie pomogło?
To zliczenie, które ktoś mógłby pójść i wykonać, co stawia rozkład a priori w
tej samej klasie co każdą inną wielkość w tym programie.

Powiedzmy, że $\val{p27.prior.real.pct}$ procent podejmowanych porównań
dotyczy prawdziwej różnicy i że test znajduje prawdziwą w
$\val{p27.power.pct}$ procentach przypadków. Jaka część porównań
raportujących $p < \val{p27.alpha}$ jest błędna?
\dotline
\nextframe
\end{fr}

\begin{fr}
\ans{$\val{p27.fdr.pct}$ procent}

Nie $\val{p27.alpha}$ z nich. Prawie połowa.

Arytmetyka to dwa zliczenia i dzielenie, czyli rzecz z
Programu~\ref{prog:F10}. Na sto porównań $\val{p27.prior.real.pct}$ jest
prawdziwych i połowa z nich zostaje znaleziona, czyli pięć prawdziwych
odkryć; dziewięćdziesiąt nie jest prawdziwych i jedno na dwadzieścia z nich
mimo to raportuje istotność, czyli cztery i pół fałszywych. Cztery i pół z
dziewięciu i pół.

\textbf{Próg rządzi fałszywymi pozytywami wśród hipotez zerowych. Nie mówi nic
o fałszywych pozytywach wśród odkryć}, a to drugie jest wielkością, która
obchodzi kogokolwiek czytającego wynik. Przesunięcie częstości bazowej
przesuwa odpowiedź, a p-wartość nie drgnie.
\end{fr}

\begin{fr}
Jeszcze jedna własność, i to ona nadaje progowi jakikolwiek sens. Jeśli oba
modele naprawdę są identyczne i przeprowadzasz porównanie, jak często $p$
wychodzi poniżej $\val{p27.alpha}$?
\dotline
\nextframe
\end{fr}

\begin{fr}
\ans{W $\val{p27.null.rate.pct}$ procentach przypadków}

Z konstrukcji, i nie jest to wada \dash{} to definicja progu widziana z drugiej
strony. Próg jest tak zbudowany, że \textbf{jedno na dwadzieścia porównań dwóch
identycznych modeli raportuje istotną różnicę}, i żaden wybór testu tego nie
usuwa.

Pomiar daje tu $\val{p27.null.rate.pct}$, a nie dokładnie $\val{p27.alpha}$, i
powód warto mieć: test na pozycjach niezgodnych jest dyskretny, więc jego
p-wartość nie może wylądować na progu, a osiągalne wartości przekraczają go
skokiem. Test dyskretny jest zatem lekko \emph{zachowawczy}, a skrypt wylicza
wszystkie $\val{p27.disc}$ wyników, zamiast zakładać odpowiedź ciągłą.

P-wartość jest więc twierdzeniem o długofalowym zachowaniu procedury, a nie o
twoim wyniku. To przydatne twierdzenie. Nie to, którego chciałeś.
\end{fr}

\begin{fr}
\begin{rigourbox}
\textbf{Czego chciałeś i gdzie to mieszka.}

Pytanie w czyjejś głowie brzmi \enquote{jakie jest prawdopodobieństwo, że B
jest lepszy od A}, i żadna p-wartość na nie nie odpowiada, bo to
prawdopodobieństwo wymaga rozkładu po twierdzeniu, a nie po danych.

Program~\ref{prog:P28} je liczy. Nałożenie rozkładu na parametr, zwarunkowanie
na ewaluacji i odczytanie $\Prob(B > A)$ to cała sekcja o rozkładzie
beta-dwumianowym w tamtym programie, a ponadto zestawia on wychodzący stamtąd
przedział z tym, który buduje ten program, bo oba odpowiadają na różne pytania
i rutynowo cytuje się je tak, jakby odpowiadały na to samo.

Ten program zatrzymuje się celowo na odczycie częstościowym: to jest to, co
raportują narzędzia stojące przed tobą, a zrozumienie, czego nie mówią, jest
warunkiem chcenia tamtej drugiej rzeczy.
\end{rigourbox}
\end{fr}

% ============================================================
\section{Czterdzieści modeli w jednym rankingu}
\label{sec:P27-many}
\index{wnioskowanie!wielokrotne porównania}
\index{ewaluacja!ranking}

\begin{fr}
Wszystko do tej pory dotyczyło jednego porównania. Ranking nie jest jednym
porównaniem.

Umieść $\val{p27.models}$ modeli na tym samym zbiorze ewaluacyjnym, a ciekawe
wielkości są parami \dash{} weźmy jednak najprostszą wersję: każdy model
porównywany z odniesieniem, czyli $\val{p27.models}$ porównań. Załóżmy, że
\textbf{żaden} z $\val{p27.models}$ w ogóle nie różni się od odniesienia.

Jaka jest szansa, że co najmniej jeden zaraportuje
$p < \val{p27.alpha}$?
\dotline
\nextframe
\end{fr}

\begin{fr}
\ans{$\val{p27.any.false.pct}$ procent}

$1 - \num{0.95}^{\val{p27.models}}$, czyli reguła dopełnienia z
Programu~\ref{prog:F10} i niezależność z Programu~\ref{prog:P23}, i nic poza
tym.

\textbf{W rankingu $\val{p27.models}$ modeli, w którym nic się nie dzieje,
prawie zawsze coś wygląda na istotne.} Próg zbudowano tak, żeby przypadek
przekraczał go raz na dwadzieścia, a czterdzieści prób to czterdzieści szans
na bycie tym jednym.

Prostacka naprawa polega na podzieleniu progu przez liczbę porównań \dash{}
$\val{p27.alpha}$ staje się $\val{p27.alpha}/\val{p27.models}$ \dash{} co
pochodzi od Bonferroniego i co robi większość dziedzin. Działa w tym sensie, że
przywraca łączną częstość. Ile kosztuje?
\dotline
\nextframe
\end{fr}

\begin{fr}
\ans{$\val{p27.bonf.cost}$ raza tyle pozycji}

Próg przesuwa mnożnik w przedziale z Programu~\ref{prog:P25} z
$\num{1.96}$ na $\val{p27.z.bonf}$, a wzór z sekcji 3 niesie ten mnożnik
\emph{do kwadratu}. Ranking, który chce czterdziestu uczciwych porównań,
potrzebuje więc $\val{p27.bonf.cost}$ raza tylu pozycji ewaluacyjnych co
jedno, na dokładkę do wszystkiego innego.

Co jest wymianą postawioną otwarcie i tak właśnie należy ją widzieć:
wielokrotne porównania nie są statystycznym drobiazgiem, tylko pozycją w
budżecie. Zadaj czterdzieści pytań jednemu zbiorowi ewaluacyjnemu, a albo
godzisz się, że jedna z odpowiedzi jest błędna, albo kupujesz prawie trzy razy
tyle pozycji.
\end{fr}

\begin{fr}
Prowadzenie wielu porównań ma drugi skutek, innego rodzaju, większy i prawie
nigdy niewspominany. Przeżywa nawet wtedy, gdy nikt nie liczy żadnej
p-wartości.

Weź $\val{p27.models}$ modeli o \textbf{dokładnie równych} prawdziwych
dokładnościach, na poziomie $\val{p27.lb.p.pct}$ procent. Zewaluuj każdy na
$\val{p27.lb.n}$ pozycjach i posortuj tabelę. Model na szczycie to ten, któremu ewaluacja poszła najlepiej, a jego
zaobserwowany wynik leży powyżej jego własnej prawdziwej dokładności.

O ile mniej więcej, w punktach? Błąd standardowy jednego wyniku to tutaj
$\val{p27.lb.se}$ punktu.
\dotline
\nextframe
\end{fr}

\begin{fr}
\ans{O jakieś $\val{p27.lb.margin}$ punktu}

Zwycięzca $\val{p27.models}$ równych zawodników leży średnio
$\val{p27.emax}$ błędu standardowego powyżej prawdy, a
$\val{p27.emax} \times \val{p27.lb.se}$ to $\val{p27.lb.margin}$.

Ta wielkość to wartość oczekiwana maksimum $\val{p27.models}$ standardowych
zmiennych normalnych i jest \textbf{scałkowana, a nie zasymulowana} \dash{} z
powodu podanego w Programie~\ref{prog:P24}, że \enquote{zgodziło się w
granicach słupka błędu} to dokładnie ten odczyt, którego taka sekcja odmawia.
Dwa przypadki z postacią zamkniętą, $m = 1$ i $m = 2$, sprawdzono wobec niej
najpierw.

\mermaidfig{p27-winner-is-noise}{Co sortowanie tabeli równych robi z jej
  szczytem}{przewaga}

W rankingu, w którym różnicę jednego punktu raportuje się jako wynik, model na
czele stawki czterdziestu wygląda więc na jakieś trzy punkty lepszy, niż jest,
\textbf{przy zerowych różnicach między modelami}. Zrobiła to selekcja, a nie
pomiar.
\end{fr}

\begin{fr}
\begin{trapbox}
\textbf{\enquote{Sprawdziliśmy trzydzieści konfiguracji i wzięliśmy tę, która
wypadła najlepiej na zbiorze ewaluacyjnym.}}

Czyli ta sama arytmetyka, powiedziana o własnej pracy zamiast o publicznej
tabeli \dash{} i to jest powód, dla którego istnieje zbiór odłożony, podany
jako liczba, a nie jako reguła, którą ktoś wręczył.

Wynik wybranej konfiguracji nie jest oceną tego, jak ona wypadnie. Jest
maksimum trzydziestu zaszumionych ocen, a maksimum $m$ rzeczy jest obciążone w
górę o wielkość rosnącą z $m$. Raportowanie go jako spodziewanej jakości to
raportowanie selekcji razem z modelem.

Naprawa nie jest statystyczna. Wybieraj na jednym zbiorze, a \emph{raportuj} z
drugiego, i liczba, którą podajesz, nie była selekcjonowana na niczym.
\end{trapbox}
\end{fr}

% ============================================================
\section{Co potrafisz teraz sprawdzić}
\label{sec:P27-check}
\index{ewaluacja!o co pytać przy opublikowanej liczbie}
\index{wnioskowanie!czego ten program nie twierdzi}

\begin{fr}
Wracamy do zdania, od którego program się zaczął, i da się już na nie
odpowiedzieć, a nie tylko w nie wątpić.

\begin{center}
\emph{Model B dostał $\val{p27.quoted.hi}$, a model A
$\val{p27.quoted.lo}$ na $\val{p27.items}$ pozycjach ewaluacyjnych.}
\end{center}

Cztery rzeczy i żadna nie wymagała biblioteki.

\textbf{Wynik nie leży na siatce.} $\val{p27.quoted.hi}$ procent z
$\val{p27.items}$ pozycji to $\val{p27.k.hi}$ pozycji, więc albo mianownikiem
nie jest $\val{p27.items}$, albo metryka nie jest dokładnością. Zapytaj które.

\textbf{Różnica to $\val{p27.gap.items}$ pozycja}, a procent jest znakomitym
sposobem, żeby to ukryć.

\textbf{Przedział wokół każdego z wyników jest ponad
$\val{p27.hw.ratio.floor}$ razy większy od różnicy}, z jednego pierwiastka.

\textbf{A test sparowany, czyli ten właściwy, daje
$\val{p27.p.exact}$} \dash{} na $\val{p27.disc}$ pozycjach niezgodnych, czyli
jedynej liczbie, o którą trzeba by poprosić. Przewaga o jeden przy
nieparzystej liczbie jest najmniejszą, jaka istnieje, więc test na pewno
zobaczy co najmniej ją.

Różnica musiałaby wynieść $\val{p27.net.needed.pts}$ punktu, zanim ten zbiór
ewaluacyjny w ogóle odróżniłby oba modele.
\end{fr}

\begin{fr}
Dwie rzeczy, których ten program nie twierdzi, i obie są warte więcej niż
kolejna technika.

\textbf{Nie zmierzył, ile opublikowanych wyników jest szumem.} Policzył, czego
trzeba, żeby jeden z nich nim nie był, na jednej przerobionej ewaluacji, a to
inne twierdzenie. Zrobienie pierwszego porządnie znaczy zebranie liczby
pozycji niezgodnych stojących za dużą liczbą opublikowanych porównań, a tych
się nie publikuje \dash{} co samo w sobie jest wynikiem i dlatego aibox z
sekcji 3 prosi o dwie liczby całkowite, a nie o lepszą statystykę.

\textbf{Ani nic z tego nie jest procedurą decyzyjną.} P-wartość poniżej progu
nie jest zgodą na wdrożenie, a powyżej nie jest dowodem równoważności:
arytmetyka z sekcji 4 mówi zwłaszcza to drugie, bo test, który potrzebuje
$\val{p27.net.needed.pts}$ punktu, żeby cokolwiek zobaczyć, nie powiedział
niczego o prawdziwej różnicy jednego.

To, co zostaje, gdy porównanie nie rozstrzyga rzeczy, na której ci zależy, jest
pytaniem o projekt, a nie o arytmetykę. Więcej pozycji, po cenie z
Programu~\ref{prog:P25}. Te same prompty, ze zniżką z sekcji 3. Albo inny
pomiar.

Te trzy rzeczy nie kosztują tyle samo i warto znać kolejność. Parowanie jest
darmowe i zwykle już dostępne. Pozycje kosztują liniowo w pieniądzach i
kwadratowo w precyzji. A inny pomiar \dash{} taki, w którym rzecz, na której
ci zależy, daje większy sygnał \dash{} jako jedyny z trzech nie ma sufitu i
jako jedyny nie stoi w niczyim budżecie.

Jest jeszcze czwarta możliwość i to ona jest najtańsza.
\dotline
\nextframe
\end{fr}

\begin{fr}
\begin{ansblock}
Albo uczciwa odpowiedź, dostępna za darmo: zaraportuj, że na tej ewaluacji obu
modeli nie dało się odróżnić, i podaj, jak duża musiałaby być różnica.
\end{ansblock}

Ten ostatni człon zmienia \enquote{brak istotnej różnicy} z niewyniku w pomiar,
i prawie nikt go nie pisze. To $\val{p27.net.needed.pts}$ punktu z sekcji 3, a
policzenie go kosztuje pętlę po \code{math.comb}.

Gdzie cię to zostawia. Umiesz przeczytać liczbę z benchmarku wobec jej własnego
mianownika; umiesz nałożyć przedział na metrykę bez wzoru i nazwać jednostkę do
losowania; umiesz zaplanować i przeprowadzić porównanie, którego ludzie
naprawdę chcą, korzystając z tego, że oba modele widziały te same prompty;
umiesz powiedzieć, co p-wartość mówi, a czego nie, nie czcząc jej ani nie
odrzucając; i umiesz wycenić rankingową porcję porównań, zanim przeczytasz jego
górny wiersz.

Program~\ref{prog:P28} bierze pytanie, które ten program musiał oddać \dash{}
prawdopodobieństwo, że B jest lepszy \dash{} i odpowiada na nie, nakładając
rozkład na to, czego nie jesteś pewien, a nie na dane, które już masz.
\end{fr}

\begin{summarybox}
\sumitem{1--3}{\result{\textbf{Czytaj wynik najpierw wobec jego własnego mianownika.}
  Wynik $k$ zaliczonych z $n$ leży na siatce o oczku $100/n$ punktu, więc
  $\val{p27.quoted.hi}$ procent na $\val{p27.items}$ pozycjach to
  $\val{p27.k.hi}$ pozycji i jest nieosiągalne. Dokładność z jednym miejscem
  po przecinku wymaga ponad tysiąca pozycji \dash{}
  $\val{p27.decimal.floor}$ \dash{} zanim to miejsce się zaokrągli}.}
\sumitem{3--5}{\result{\textbf{Opublikowana różnica to $\val{p27.gap.items}$ pozycja.}
  $\val{p27.gap.pts}$ punktu z $\val{p27.items}$ pozycji, wobec przedziału
  $\pm\val{p27.hw.200}$ punktu \dash{} ponad
  $\val{p27.hw.ratio.floor}$ razy większego od różnicy}.}
\sumitem{6}{\result{Porównywanie dwóch osobnych przedziałów po tym, czy się nakładają,
  nie jest testem. Buduj przedział na różnicy, czyli na wielkości, o którą
  pytasz}.}
\sumitem{7--9}{\result{\textbf{Bootstrap} losuje z próbki ze zwracaniem i nie
  potrzebuje wzoru dla żadnej metryki. Na proporcji jego rozkład jest
  \emph{dokładnie} dwumianowy, więc zgadza się ze wzorem z
  Programu~\ref{prog:P25} co do jednej pozycji z konstrukcji, a nie
  przypadkiem}.}
\sumitem{10--11}{\result{Jego zastosowaniem jest przypadek bez wzoru \dash{} mediana,
  której przedział $\intcc{\val{p27.med.lo}}{\val{p27.med.hi}}$ ma końce
  będące pomiarami, które ktoś wykonał. Nie wymyśli pozycji, których nie
  ewaluowałeś, i najgorzej radzi sobie ze skrajnościami}.}
\sumitem{12--14}{\result{\textbf{Jednostka losowania to decyzja, której nic nie
  sprawdza.} $\val{p27.clust.n}$ wierszy od $\val{p27.clust.users}$
  użytkowników to nie $\val{p27.clust.n}$ obserwacji, a bootstrap po wierszach
  podaje przedział za wąski o $\val{p27.clust.factor}$. Więcej losowań nigdy
  nie zwęża przedziału; robią to tylko pozycje}.}
\sumitem{15--17}{\result{\textbf{Parowanie mnoży potrzebne pozycje przez
  $1-\rho$}, dokładnie, przy każdej dokładności i każdej różnicy.
  Program~\ref{prog:P25} podaje $\val{p27.n.rho00}$, co staje się
  $\val{p27.n.rho90}$ przy $\rho = \num{0.9}$ \dash{} rząd wielkości, z wyrazu
  kowariancji w $\Var(X-Y)$}.}
\sumitem{18--19}{\result{\textbf{Różnicę niosą tylko pozycje niezgodne.} Te
  $\val{p27.concord}$ pozycji, na których oba modele odpowiedziały tak samo,
  skracają się dokładnie, a porównanie dotyczy w całości
  $\val{p27.disc}$ pozostałych}.}
\sumitem{19--20}{\result{Przy \enquote{równie dobre} te pozycje są uczciwymi rzutami
  monetą, więc testem jest symbol Newtona z Programu~\ref{prog:P12} na
  liczbach całkowitych \dash{} dokładny, niezależny od maszyny i niepotrzebujący
  wariancji. \textbf{Przerobiona różnica daje $\val{p27.p.exact}$}: nieparzysta
  liczba nie wypada po równo, więc zaobserwowana przewaga
  $\val{p27.net}$ jest najmniejszą dostępną, a wynik co najmniej tak
  przechylony jest pewny}.}
\sumitem{21}{\result{\textbf{Na tej ewaluacji różnica poniżej
  $\val{p27.net.needed.pts}$ punktu nie jest wynikiem}, znalezione przez
  szukanie. Opublikowana różnica to $\val{p27.gap.pts}$ punktu}.}
\sumitem{23--25}{\result{\textbf{P-wartość to
  $\Prob(\text{dane tak skrajne} \mid \text{brak różnicy})$}, a czytanie jej
  wstecz to odwrócenie z Programu~\ref{prog:P23} \dash{} ten sam błąd co przy
  częstości bazowej}.}
\sumitem{26--27}{\result{Mając rozkład a priori da się ją odwrócić, a odpowiedź
  niepokoi: przy $\val{p27.prior.real.pct}$ procentach porównań prawdziwych i
  mocy $\val{p27.power.pct}$ procent $\val{p27.fdr.pct}$ procent
  \enquote{istotnych} odkryć jest fałszywych}.}
\sumitem{28--29}{\result{\textbf{Jedno na dwadzieścia porównań identycznych modeli
  raportuje różnicę}, z konstrukcji. Tutaj $\val{p27.null.rate.pct}$ procent,
  bo test dyskretny przekracza próg skokiem, zamiast na nim lądować}.}
\sumitem{31--33}{\result{\textbf{$\val{p27.models}$ porównań robi z tego
  $\val{p27.any.false.pct}$ procent.} Podzielenie progu to przywraca i
  kosztuje $\val{p27.bonf.cost}$ raza tyle pozycji \dash{} wielokrotne
  porównania są pozycją w budżecie, a nie drobiazgiem}.}
\sumitem{34--35}{\result{\textbf{A sortowanie tabeli obciąża jej górny wiersz.}
  Zwycięzca $\val{p27.models}$ \emph{równych} modeli leży $\val{p27.emax}$
  błędu standardowego powyżej własnej prawdy, czyli $\val{p27.lb.margin}$
  punktu na $\val{p27.lb.n}$ pozycjach. Selekcja, nie pomiar \dash{} i ta sama
  arytmetyka jest powodem, dla którego istnieje zbiór odłożony}.}
\sumitem{37--39}{\result{To, co zostaje, gdy porównanie nie rozstrzyga rzeczy, na
  której ci zależy, jest pytaniem o projekt. \textbf{A \enquote{brak istotnej
  różnicy} staje się pomiarem w chwili, gdy podasz, jak duża musiałaby być
  różnica.}}}
\end{summarybox}

\canyou

\begin{testexercises}
\item Model raportuje $\num{83.3}$ procent na $60$ pozycjach ewaluacyjnych.
  Czy to osiągalne i jakie są dwa najbliższe osiągalne wyniki?
  \answerto{Tak \dash{} ale nie z powodu, dla którego ramka 2 odrzuciła
  $\num{71.4}$. $50/60$ to $\num{83.33}\ldots$ procent, co drukuje się jako
  $\num{83.3}$, więc podana liczba jest zaokrągleniem osiągalnego wyniku; ani
  $\num{71.0}$, ani $\num{71.5}$ nie zaokrągla się do $\num{71.4}$ i dlatego
  tamta nie była. Siatka ma tu oczko $100/60 = \num{1.67}$ punktu, więc
  sąsiadami są $49/60 = \num{81.7}$ i $51/60 = \num{85.0}$ procent \dash{}
  tak gruboziarnisty jest benchmark o $60$ pozycjach.}
\item Dwa modele zewaluowano na $500$ wspólnych promptach. Nie zgadzają się na
  $80$ z nich, a z tych $46$ idzie do B i $34$ do A. Ile wynosi przewaga, w
  pozycjach i w punktach?
  \answerto{Przewaga to $46 - 34 = 12$ pozycji, czyli
  $12/500 = \num{2.4}$ punktu. Zwróć uwagę, że obie dokładności różnią się o
  $\num{2.4}$ punktu niezależnie od tego, ile pozycji oba zaliczyły \dash{}
  zgodne się skracają.}
\item Ewaluacja ma $600$ wierszy pochodzących z $50$ rozmów. Kolega robi
  bootstrap po wierszach i podaje przedział. W którą stronę jest on błędny i
  mniej więcej o jaki czynnik, jeśli wiersze wewnątrz rozmowy zgadzają się
  całkowicie?
  \answerto{Za wąski, nawet o $\sqrt{600/50} = \sqrt{12} \approx
  \num{3.5}$. Bootstrap sądzi, że ma $600$ niezależnych obserwacji, a w
  najgorszym razie ma $50$.}
\item Zespół przeprowadza $20$ porównań z odniesieniem i raportuje dwa, które
  wróciły z $p < \num{0.05}$. Gdyby żaden z $20$ systemów nie różnił się od
  odniesienia, ilu takich raportów byś się spodziewał?
  \answerto{$20 \times \num{0.05} = 1$. Znalezienie dwóch nie jest dowodem
  niczego; oczekiwana liczba przy zerowej różnicy wynosi jeden, a
  $\Prob(\text{co najmniej jeden}) = 1 - \num{0.95}^{20} = \num{0.64}$.}
\item Twoje porównanie potrzebuje $\val{p27.n.rho00}$ pozycji na model, jeśli
  obie ewaluacje są niezależne. Masz $\val{p27.n.rho80}$ promptów i możesz
  puścić oba modele na wszystkich. Jakiej korelacji wyników obu modeli na
  poszczególnych pozycjach potrzebowałbyś, żeby to wystarczyło?
  \answerto{$\val{p27.n.rho00} \times (1 - \rho) \le \val{p27.n.rho80}$ daje
  $\rho \ge \num{0.8}$ \dash{} czyli to, co mówi tabela z sekcji 3, czytana
  wstecz.}
\end{testexercises}

\begin{furtherproblems}
\item Benchmark raportuje $\num{62.48}$ procent. Jaka jest najmniejsza liczba
  pozycji, na której dałoby się w ogóle policzyć tyle miejsc po przecinku?
  \answerto{Ten sam warunek co w ramce 3, o jedno miejsce dalej: siatka musi
  być gęstsza niż setna punktu, więc $1/n < \num{0.0001}$ i $n > 10\,000$
  \dash{} czyli $\val{p27.decimal.floor.two}$ pozycji. Prawie żaden benchmark
  nie ma tylu, więc dwa miejsca po przecinku w publicznym rankingu są
  ozdobą.}
\item Pokaż, że rozkład bootstrapowy \emph{średniej} z $n$ wartości ma
  wariancję $(1 - 1/n)$ razy zwykłe $s^{2}/n$, i powiedz, dlaczego nie ma to
  znaczenia przy $n = \val{p27.items}$.
  \answerto{Losowanie idzie z rozkładu empirycznego, którego wariancja jest tą
  z mianownikiem $n$ \dash{} Program~\ref{prog:P26} nazywa ten czynnik
  $\frac{n-1}{n}$. Przy $n = \val{p27.items}$ to różnica
  $\num{0.5}$ procent w wariancji i ćwierć procent w przedziale, czyli
  głęboko wewnątrz siatki zmierzonej w sekcji 1.}
\item Dwa modele porównano na $n$ wspólnych pozycjach i nie zgadzają się na
  ułamku $d$ z nich. Zapisz liczbę pozycji potrzebną do rozstrzygnięcia
  różnicy $\delta$ przez $d$ zamiast przez $\rho$ i powiedz, którą z tych
  dwóch wielkości harness ewaluacyjny mógłby faktycznie raportować.
  \answerto{Ułamek pozycji niezgodnych jest tym, co harness umie zliczyć, i
  jest wielkością, której test dokładny używa wprost: przewaga musi sięgnąć
  około $\num{1.96}\sqrt{dn}$ pozycji, żeby uczciwa moneta jej nie
  produkowała \dash{} zaokrąglone w górę do kolejnej wartości o tej samej
  parzystości co $dn$, i z luzem: przy $dn = \val{p27.disc}$ daje to $11$,
  gdzie dokładne przeszukanie daje $\val{p27.net.needed}$, bo test dyskretny
  przeskakuje swój próg, zamiast na nim wylądować. $\rho$ trzeba oszacować;
  $d$ to dwa liczniki w pętli, i
  dlatego aibox z sekcji 3 prosi właśnie o nie.}
\item Dwa czołowe modele rankingu dzieli $\num{0.3}$ punktu na
  $\val{p27.lb.n}$ pozycjach. Bez żadnych dalszych informacji, co
  najpożyteczniejszego da się powiedzieć?
  \answerto{Że dzieli je $3$ pozycje, że błąd standardowy każdego z wyników to
  $\val{p27.lb.se}$ punktu i że przy $\val{p27.models}$ modelach w tabeli
  górny wiersz leży spodziewanie o jakieś $\val{p27.lb.margin}$ punktu powyżej
  własnej prawdziwej dokładności z samej selekcji \dash{} czyli ponad
  dziesięć razy więcej niż raportowana różnica.}
\item Program~\ref{prog:P26} pokazał, że celowo obciążony estymator potrafi
  pobić nieobciążony. Czy ta sama wymiana dotyczy przedziału ufności i jakie
  byłyby obie osie?
  \answerto{Tak, a osiami są \emph{pokrycie} i \emph{szerokość}. Przedział,
  który zawsze wynosi $\intcc{0}{100}$, ma doskonałe pokrycie i zero pożytku;
  ten, który zawsze jest punktem, ma zerową szerokość i zero pokrycia.
  Przedział konwencjonalny ustala pokrycie na $95$ procentach i podaje
  szerokość, jaką to kosztuje, co jest jednym wyborem na tej wymianie, a nie
  jedynym.}
\end{furtherproblems}
